% Part: computability
% Chapter: recursive-functions
% Section: sequences

\documentclass[../../../include/open-logic-section]{subfiles}

\begin{document}

\olfileid{cmp}{rec}{seq}
\olsection{క్రమాలు}

ఆదిమ పునరావృత్త ప్రమేయాల సమితి విశేషంగా దృఢమైనది. అయితే
\emph{క్రమాలను} నిర్వహించడానికి తగిన పద్ధతిని అభివృద్ధి చేసిన తరువాత
ఇంకా ఎక్కువ చేయగలం. సహజ సంఖ్యల పరిమిత క్రమాలను సహజ సంఖ్యలతో కింది
విధంగా గుర్తిస్తాం: క్రమం $\langle a_0,a_1,a_2,\dots,a_k\rangle$కు
అనుగుణమైన సంఖ్య
\[
p_0^{a_0+1} \cdot p_1^{a_1+1} \cdot p_2^{a_2+1} \cdot \dots \cdot
p_k^{a_k+1}.
\]
ఉదాహరణకు, $\langle2,7,3\rangle$, $\langle2,7,3,0,0\rangle$లకు
వేర్వేరు సంఖ్యా సంకేతాలు ఉండేలా ఘాతాలకు ఒకటి కలుపుతాం. ఖాళీ క్రమానికి
$0$,~$1$ రెండింటినీ సంకేతాలుగా తీసుకోవచ్చు; నిర్దిష్టత కోసం
$\emptyseq$ అనేది~$0$ను సూచించనివ్వండి.

క్రమాల ఈ సంకేతీకరణ పనిచేయడానికి కారణం అంకగణిత ప్రాథమిక సిద్ధాంతం:
ప్రతి సహజ సంఖ్య $n\ge2$ను $a_k\ge1$ అయ్యేలా
\[
n = p_0^{a_0} \cdot p_1^{a_1} \cdot \dots \cdot p_k^{a_k}
\]
రూపంలో ఒకే ఒక్క విధంగా రాయవచ్చు. అందువల్ల
$(a_0,\dots,a_k)\longmapsto\tuple{a_0,\dots,a_k}$ అనే ప్రతిబింబం
ఒకటి-ఒకటి: వేర్వేరు క్రమాలు వేర్వేరు సంఖ్యలకు వెళ్తాయి; ప్రతి
సంఖ్యకు అత్యధికంగా ఒక క్రమమే అనుగుణమవుతుంది.
\sourcecorrection{OLTECRREM-006}{ఆంగ్ల మూలంలోని సంకేతీకరణ ప్రతిబింబం $\tuple{}(a_0,\dots,a_k)=\tuple{a_0,\dots,a_k}$గా వికృతమైంది. వాక్యం ప్రకటించే ఒకటి-ఒకటి ప్రతిబింబాన్ని $(a_0,\dots,a_k)\longmapsto\tuple{a_0,\dots,a_k}$గా స్పష్టం చేశాం.}

ఒక క్రమపు పొడవును కనుగొనడం, దాని $i$వ మూలకాన్ని కనుగొనడం, క్రమం
చివర ఒక మూలకాన్ని చేర్చడం, రెండు క్రమాలను అనుసంధానించడం అనే క్రియలన్నీ
ఆదిమ పునరావృత్తాలని ఇప్పుడు చూపుతాం.

\begin{prop}
  క్రమం $s$ యొక్క పొడవును తిరిగి ఇచ్చే ప్రమేయం $\len{s}$ ఆదిమ
  పునరావృత్త ప్రమేయం.
\end{prop}

\begin{proof}
  సంబంధం $R(i,s)$ను ఇలా నిర్వచిద్దాం:
  \[
  R(i, s) \text{ అప్పుడే }
  p_i \mid s \land p_{i+1} \nmid s.
  \]
  $R$ స్పష్టంగా ఆదిమ పునరావృత్త సంబంధం. $s$ ఒక ఖాళీకాని క్రమపు
  సంకేతమైనప్పుడల్లా, అంటే
  \[
  s = p_0^{a_0+1} \cdot \dots \cdot p_{k}^{a_{k}+1},
  \]
  అయినప్పుడు, $s$ను భాగించే ప్రధాన సంఖ్యల్లో $p_i$ అతి పెద్దదై
  $p_i \mid s$ అయినప్పుడే,
  అంటే $i=k$ అయినప్పుడే, $R(i,s)$ నెరవేరుతుంది. కాబట్టి $p_i$ అనేది
  $s$ను భాగించే అతి పెద్ద ప్రధాన సంఖ్య అయినప్పుడే $s$ పొడవు $i+1$;
  అందువల్ల
  \[
  \len{s} =
  \begin{cases}
    0 & \text{$s=0$ లేదా $s=1$ అయితే} \\
    1 + \bmin{i < s}{R(i, s)} & \text{ఇతర సందర్భాల్లో.}
  \end{cases}
  \]
  అని తీసుకోవచ్చు. $s$ ఒక క్రమపు సంకేతమైతే $R(i,s)$ను నెరవేర్చే
  $i$ ఒక్కటే ఉంటుంది; అటువంటి $i$ ఉంటే అది $s$కన్నా చిన్నదే. కనుక
  ఇక్కడ పరిమిత
  కనిష్ఠీకరణను ఉపయోగించవచ్చు.
\end{proof}

\begin{prop}
  క్రమం $s$ చివర $a$ను చేర్చిన ఫలితాన్ని తిరిగి ఇచ్చే ప్రమేయం
  $\fn{append}(s,a)$ ఆదిమ పునరావృత్త ప్రమేయం.
\end{prop}

\begin{proof}
  $\fn{append}$ను ఇలా నిర్వచించవచ్చు:
  \[
  \fn{append}(s,a) =
  \begin{cases}
    2^{a+1} & \text{$s=0$ లేదా $s=1$ అయితే} \\
    s \cdot p_{\len{s}}^{a+1} & \text{ఇతర సందర్భాల్లో.}
  \end{cases}
  \]
\end{proof}

\begin{prop}
  $s$ యొక్క $i$వ మూలకాన్ని (మొదటి మూలకాన్ని $0$వది అంటాం), లేదా
  $i$ అనేది $s$ పొడవుకన్నా ఎక్కువ లేదా సమానమైతే $0$ను తిరిగి ఇచ్చే
  ప్రమేయం $\fn{element}(s,i)$ ఆదిమ పునరావృత్త ప్రమేయం.
\end{prop}

\begin{proof}
$p_i^{a+1}$ అనేది $s$ను భాగించే $p_i$ యొక్క అతి పెద్ద ఘాతం
అయినప్పుడే $a$ అనేది $s$ యొక్క $i$వ మూలకం; అంటే
$p_i^{a+1}\mid s$, కానీ $p_i^{a+2}\nmid s$. కాబట్టి:
  \[
  \fn{element}(s,i) =
  \begin{cases}
    0 & \mbox{$i\geq\len{s}$ అయితే} \\
    \bmin{a < s}{(p_i^{a+2} \nmid s)} & \text{ఇతర సందర్భాల్లో.}
  \end{cases}
  \]
\end{proof}

పైన నిర్వచించిన ప్రమేయాల అధికారిక పేర్లకు బదులుగా మరింత క్లుప్తమైన
సంకేతనాన్ని ప్రవేశపెడతాం. $\fn{element}(s,i)$కు బదులు $(s)_i$ని,
కింది వ్యక్తీకరణకు సంక్షిప్తంగా $\tuple{s_0,\dots,s_k}$ను వాడతాం:
\[
\fn{append}(\fn{append}(\dots \fn{append}(\emptyseq,s_0)
\dots),s_k).
\]
$s$ పొడవు~$k$ అయితే దాని మూలకాలు $(s)_0$, \dots,~$(s)_{k-1}$ అని
గమనించండి.

\begin{prop}
రెండు క్రమాలను అనుసంధానించే ప్రమేయం $\fn{concat}(s,t)$ ఆదిమ
పునరావృత్త ప్రమేయం.
\end{prop}

\begin{proof}
  మనకు కింది ధర్మం గల $\fn{concat}$ ప్రమేయం కావాలి:
  \[
    \fn{concat}(\tuple{a_0, \dots, a_k}, \tuple{b_0, \dots, b_l}) =
    \tuple{a_0, \dots, a_k, b_0, \dots, b_l}.
  \]
  $t$లోని మొదటి $n$ మూలకాలను $s$కు అనుసంధానించే ``సహాయక'' ప్రమేయం
  $\fn{hconcat}(s,t,n)$ను వాడదాం. దాన్ని ఆదిమ పునరావృత్తితో ఇలా
  నిర్వచించవచ్చు:
  \begin{align*}
    \fn{hconcat}(s,t,0) & = s\\
    \fn{hconcat}(s,t,n+1) & = \fn{append}(\fn{hconcat}(s,t,n),(t)_n)
    \intertext{అప్పుడు $\fn{concat}$ను ఇలా నిర్వచించవచ్చు:}
    \fn{concat}(s,t) & = \fn{hconcat}(s,t,\len{t}).
  \end{align*}
\end{proof}

$\fn{concat}(s,t)$కు బదులుగా $s\concat t$ అని రాస్తాం.

ఒక క్రమపు సంఖ్యా సంకేతాన్ని దాని పొడవు, దాని అతి పెద్ద మూలకం పరంగా
అవధి చేయగలగడం మనకు ఉపయోగపడుతుంది. $s$ పొడవు~$k$ అనీ, దానిలోని ప్రతి
మూలకం ఏదో సంఖ్య~$x$కన్నా చిన్నది లేదా సమానం అనీ అనుకుందాం. అప్పుడు
$s$కు అత్యధికంగా $k$ ప్రధాన కారణాంకాలు ఉంటాయి; ప్రతి కారణాంకం
అత్యధికంగా $p_{k-1}$, $s$ యొక్క ప్రధాన కారణాంక విభజనలో ప్రతి ఘాతం అత్యధికంగా
$x+1$. ఖాళీ క్రమానికి అవధిని 1గా తీసుకుంటాం; పొడవు ధనాత్మకమైనప్పుడు
\[
\fn{sequenceBound}(x,k) = p_{k-1}^{k \cdot (x+1)},
\]
అని నిర్వచిస్తే, పైన వర్ణించిన క్రమం~$s$ యొక్క సంఖ్యా సంకేతం
అత్యధికంగా $\fn{sequenceBound}(x,k)$.

క్రమాలపై ఇటువంటి అవధి ఉండటం వల్ల పరిమిత అన్వేషణను ఉపయోగించి కొత్త
ప్రమేయాలను నిర్వచించవచ్చు. ఉదాహరణకు, $\fn{concat}$ను పరిమిత
అన్వేషణతో నిర్వచించవచ్చు. మనం వెతుకుతున్న వస్తువుకు (అనుసంధానించిన
క్రమపు సంఖ్యకు) ఒక ఆదిమ పునరావృత్త \emph{నిర్దేశనను}, ఎంతవరకు
వెతకాలో ఒక అవధిని రాస్తే చాలు. కిందిది పనిచేస్తుంది:
\begin{align*}
  \fn{concat}(s,t) = {} & \bmin{v \leq \fn{sequenceBound}(s+t,\len{s} +
    \len{t})}{} \\
  & \quad(\len{v} = \len{s} + \len{t} \land {}\\
    & \qquad \bforall{i < \len{s}}{((v)_i = (s)_i) \land {} \\
      & \qquad \bforall{j < \len{t}}{((v)_{\len{s}+j} = (t)_j)})}
\end{align*}
\sourcecorrection{OLTECRREM-007}{ఆంగ్ల మూలంలోని $\fn{sequenceBound}$ సూత్రం $k=0$ వద్ద $p_{-1}$ను కోరింది; ఖాళీ క్రమానికి అవధి $1$ను చేర్చాం. అలాగే సంకేతం ఆ అవధికి సమానం కావచ్చునప్పటికీ మూలం $v<\fn{sequenceBound}$ వరకే వెతికింది; నిరూపిత ``అత్యధికంగా'' అవధికి అనుగుణంగా $v\leq\fn{sequenceBound}$గా సరిచేశాం.}

\begin{prob}
కింది ధర్మం గల ఆదిమ పునరావృత్త ప్రమేయం~$\fn{sconcat}(s)$ ఉందని
చూపండి:
\[
\fn{sconcat}(\tuple{s_0, \dots, s_k}) = s_0 \concat \dots \concat s_k.
\]
\end{prob}

\begin{prob}
కింది ధర్మాలు గల ఆదిమ పునరావృత్త ప్రమేయం~$\fn{tail}(s)$ ఉందని
చూపండి:
\begin{align*}
  \fn{tail}(\emptyseq) & = 0 \text{ మరియు}\\
  \fn{tail}(\tuple{s_0, \dots, s_{k}}) & = \tuple{s_1, \dots, s_{k}}.
\end{align*}
\end{prob}

\begin{prop}
  \ollabel{prop:subseq}
  $s$ యొక్క $i$వ మూలకం వద్ద మొదలై పొడవు~$n$ ఉన్న ఉపక్రమాన్ని తిరిగి
  ఇచ్చే ప్రమేయం $\fn{subseq}(s,i,n)$ ఆదిమ పునరావృత్త ప్రమేయం.
\end{prop}

\begin{proof}
  అభ్యాసం.
\end{proof}

\begin{prob}
  \olref[cmp][rec][seq]{prop:subseq}ను నిరూపించండి.
\end{prob}

\end{document}
